% !PRINTED 61
\weberpage{61}

Wollten wir $l_1(u)$, $l_2(u)$ als Konstanten annehmen, so wäre $T(u)$ eine ganze doppeltperiodische Funktion der zweiten Art und folglich nach \S\,16, 6. eine Exponentialfunktion, aus der sich keine doppeltperiodischen Funktionen bilden lassen. Die einfachste Annahme, die uns hiernach bleibt, ist die, daß $l_1(u)$, $l_2(u)$ \gesperrt{lineare Funktionen} von $u$ sind, worunter als Spezialfall auch der enthalten ist, daß eine der beiden Funktionen $l_1$, $l_2$ konstant ist.

Wir stellen also die folgende Definition auf:

1. \gesperrt{Eine $T$-Funktion, $T(u)$, ist eine ganze Funktion von $u$, die den folgenden beiden Gleichungen genügt:}

\weq{6}{\begin{aligned}
T(u+\omega_1) &= e^{-\pi i [a_1(2u+\omega_1)+b_1]}\,T(u),\\
T(u+\omega_2) &= e^{-\pi i [a_2(2u+\omega_2)+b_2]}\,T(u).
\end{aligned}}

Hierin sind $a_1$, $b_1$, $a_2$, $b_2$ Konstanten, d.\,h. von $u$ unabhängige Größen. Die Perioden $\omega_1$, $\omega_2$ gelten hier durchweg als ein für allemal gegebene Konstanten, deren Verhältnis $\omega_2:\omega_1$ einen positiven imaginären Teil hat.

Die Faktoren

\[
\begin{aligned}
e_1 &= e_1(u) = e^{-\pi i [a_1(2u+\omega_1)+b_1]},\\
e_2 &= e_2(u) = e^{-\pi i [a_2(2a+\omega_2)+b_2]}
\end{aligned}
\]
% [sic] The source prints 2a in the second exponent; mathematically expected 2u.

heißen die beiden \gesperrt{Periodizitätsfaktoren} der Funktion $T$.

Die Periodizitätsfaktoren ändern sich nicht, wenn $b_1$, $b_2$ um \gesperrt{gerade ganze Zahlen} verändert werden. Daß es Funktionen dieser Art wirklich gibt, wird sich erst im weiteren Verlaufe der Untersuchung ergeben.

2. \gesperrt{Wenn die Funktion $T(u)$ in irgend einem Punkte verschwindet, so verschwindet sie auch in allen kongruenten Punkten. Verschwindet sie in $m$ Punkten des Periodenparallelogramms, so heißt sie von der $m$ten Ordnung.}

Aus dieser Definition ergiebt sich sofort:

3. \gesperrt{Durch Multiplikation zweier $T$-Funktionen $T$, $T'$ der Ordnung $m$, $m'$ erhält man eine $T$-Funktion der Ordnung $m+m'$. Die Periodizitätsfaktoren des Produktes $T$, $T'$ sind die Produkte der Periodizitätsfaktoren von $T$ und $T'$.}
